\documentclass[10pt,a4paper]{amsart}
\usepackage[margin=19mm]{geometry}
\usepackage{amsmath,amssymb,mathtools}
\usepackage[scr=rsfs,cal=euler]{mathalfa}
\usepackage{xfrac}
\usepackage[unicode,hidelinks,bookmarks=false]{hyperref}
\newcommand{\ie}{i.e.\ }
\newcommand{\cf}{cf.\ }
\newcommand{\resp}{respectively\ }
\newcommand{\Subsection}{\subsection}
\providecommand{\nobreakdash}{\nobreak\hbox{-}}
\setlength{\emergencystretch}{2em}
\multlinegap0pt
\usepackage{bm}
\usepackage{bbm}
\usepackage{dsfont}
\usepackage{xspace}
\usepackage{cancel}
\usepackage{mathtools}
\usepackage{cleveref}
\usepackage{amsthm}
\makeatletter
\@ifundefined{theorem}{\newtheorem{theorem}{Theorem}}{}
\makeatother
\title[Cyclic codes over a commutative non-unitary ring of order 4]{Cyclic codes over a commutative non-unitary ring of order 4}
\author{Marvin Olavides, Jon-Lark Kim}
\date{Source: arXiv:2607.11965v5 (CC BY 4.0)}
\begin{document}
\maketitle

\noindent\emph{Source and licence.} Cyclic codes over a commutative non-unitary ring of order 4. Version arXiv:2607.11965v5 (CC BY 4.0), distributed under the Creative Commons Attribution 4.0 International licence, as recorded in the arXiv licence metadata for this exact version and shown on the abstract page https://arxiv.org/abs/2607.11965v5. Full rights notice: licenses/arxiv-2607.11965-CC-BY-4.0.txt.\par\smallskip

\noindent\textit{Editorial scope.} Counted unit: the Theorem whose source label is \texttt{thm:dual-phi} in the article (source0.tex lines 1188--1208, source label \texttt{thm:dual-phi}), delivered together with the article's own complete original proof of it (source0.tex lines 1210--1319, one \texttt{proof} environment of 110 lines). No mathematics is written, repaired or re-derived: statement and proof are the article's own text line for line. Required context retained uncounted: thm:GrayImageGeneral (lines 1080--1098). Every definition, display and label of those lines is retained unchanged. Omitted from this package and not redistributed: 1--1079, 1099--1187, 1209--1209, 1320--4208. Nothing inside a retained excerpt is omitted or altered: every retained body is delivered exactly as the article's lines are written, so no omission receipt of any kind accompanies this package. Citations: the retained excerpts carry no citation command at all, so no bibliography entry of the article is transcribed and no reference is redistributed. Reference adaptation: none was needed and none was made; every reference inside the delivered unit resolves inside it, so no unresolved reference or citation is produced. Printed slips kept literal: no printed word, symbol, hypothesis or number of the counted statement or of its proof was altered. Layout and class: the article's own class is \texttt{article} with packages including cite, amsmath, amssymb, amsfonts, mathtools, algorithmic, graphicx, textcomp, xcolor, amsthm, mathabx, xurl, footmisc, hyperref; the wrapper is the lane's pinned \texttt{amsart} editorial wrapper at 10pt on A4, which restates the theorem-like environments the retained text needs and adds the article's own macro definitions that the retained text needs (none), with extra wrapper packages (bm, bbm, dsfont, xspace, cancel, mathtools, cleveref). The delivered numbering is the unit's own. Assets: no retained span contains a figure environment, a graphics-inclusion command, a PDF-page inclusion, a table or a TikZ picture; no image file is copied into this package, the package declares no asset dependency and no image file is redistributed with it. Licence: version arXiv:2607.11965v5 of this article is distributed under the Creative Commons Attribution 4.0 International licence, as recorded in the arXiv licence metadata for this exact version and shown on the abstract page https://arxiv.org/abs/2607.11965v5. A full rights notice is delivered at licenses/arxiv-2607.11965-CC-BY-4.0.txt. Characters: every retained excerpt is pure ASCII, so the delivered engine, XeLaTeX with its default Latin Modern text font, reports no missing character.
\begin{theorem}   \label{thm:GrayImageGeneral}
Let $\mathcal{C} \subseteq I_2^n$ be an $I_2$-code. For each
$\mathbf{r} \in \operatorname{res}(\mathcal{C})$, choose a representative
$\mathbf{s}_{\mathbf{r}} \in \mathbb{F}_2^n$ of the coset
$\tau (\mathbf{r}) \in \mathbb{F}_2^n / \operatorname{tor}(\mathcal{C})$. Then the Gray image of
$\mathcal{C}$ is given by
$
\Phi(\mathcal{C})
=
\{
(\mathbf{r} \, | \,  \mathbf{r} + \mathbf{s}_{\mathbf{r}} + \mathbf{t})
\mathrel{:}
\mathbf{r} \in \operatorname{res}(\mathcal{C}), \
\mathbf{t} \in \operatorname{tor}(\mathcal{C})
\}
%\subseteq \mathbb{F}_2^{2n}.
$.
Moreover, this expression is independent of the choice of representatives $\mathbf{s}_{\mathbf{r}}$.
\end{theorem}

\begin{theorem}   \label{thm:dual-phi}
Let $\mathcal{C} \subseteq I_2^n$ be an $I_2$-code. For each
$\mathbf{r} \in \operatorname{res}(\mathcal{C})$, choose a representative
$\mathbf{s}_{\mathbf{r}} \in \mathbb{F}_2^n$ of the coset
$\tau (\mathbf{r}) \in \mathbb{F}_2^n / \operatorname{tor}(\mathcal{C})$. Then the dual of the Gray image of $\mathcal{C}$ is given by
$
\Phi(\mathcal{C})^\perp
=
\{
(\mathbf{u} \, | \, \mathbf{v})    \mathrel{:}
\mathbf{v} \in \operatorname{tor} ( \mathcal{C} )^\perp
,  \,
(\mathbf{u} + \mathbf{v}) \cdot \mathbf{r}
=
\mathbf{v} \cdot \mathbf{s}_{\mathbf{r}}  \
\forall \, \mathbf{r} \in \operatorname{res} ( \mathcal{C} )
\}
%\subseteq \mathbb{F}_2^{2n}.
$.
Moreover, this expression is independent of the choice of representatives $\mathbf{s}_{\mathbf{r}}$.
\end{theorem}

\begin{proof}
From \Cref{thm:GrayImageGeneral}, we have
\begin{align*}
\Phi(\mathcal{C})
=
\{
(\mathbf{r} \, |  \,  \mathbf{r} + \mathbf{s}_{\mathbf{r}} + \mathbf{t})
\mathrel{:}
\mathbf{r} \in \operatorname{res}(\mathcal{C})  ,  \
\mathbf{t} \in \operatorname{tor}(\mathcal{C})
\}  .
\end{align*}
Let $(\mathbf{u} \, | \, \mathbf{v}) \in \mathbb{F}_2^{2n}$. Then
$(\mathbf{u} \, | \, \mathbf{v}) \in \Phi(\mathcal{C})^\perp$ if and only if
$
(\mathbf{u} \, | \,  \mathbf{v})
\cdot
(\mathbf{r} \, |  \,  \mathbf{r} + \mathbf{s}_{\mathbf{r}} + \mathbf{t})
= 0
$
for all $\mathbf{r} \in \operatorname{res}(\mathcal{C})$ and $\mathbf{t} \in \operatorname{tor}(\mathcal{C})$. Computing the inner product
gives
\begin{align*}
0
=
(\mathbf{u} \, | \,  \mathbf{v})
\cdot
(\mathbf{r} \, |  \,  \mathbf{r} + \mathbf{s}_{\mathbf{r}} + \mathbf{t})    %\\
=
\mathbf{u} \cdot \mathbf{r}
+
\mathbf{v} \cdot (\mathbf{r} + \mathbf{s}_{\mathbf{r}} + \mathbf{t})   %\\
=
(\mathbf{u} + \mathbf{v}) \cdot \mathbf{r}
+
\mathbf{v} \cdot \mathbf{s}_{\mathbf{r}}
+
\mathbf{v} \cdot \mathbf{t}
\end{align*}
for all $\mathbf{r} \in \operatorname{res}(\mathcal{C})$ and $\mathbf{t} \in \operatorname{tor}(\mathcal{C})$.

First, fix $\mathbf{r} \in \operatorname{res}(\mathcal{C})$ and let $\mathbf{t}$ vary over $\operatorname{tor}(\mathcal{C})$. The above
expression must vanish for all $\mathbf{t} \in \operatorname{tor}(\mathcal{C})$, so we must have
$\mathbf{v} \cdot \mathbf{t} = 0$ for all $\mathbf{t} \in \operatorname{tor}(\mathcal{C})$, that is,
$
\mathbf{v} \in \operatorname{tor}(\mathcal{C})^\perp.
$
With this condition imposed, the orthogonality condition simplifies to
$
(\mathbf{u} + \mathbf{v}) \cdot \mathbf{r}
+
\mathbf{v} \cdot \mathbf{s}_{\mathbf{r}}
=
0
$
for all $\mathbf{r} \in \operatorname{res}(\mathcal{C})$,
or equivalently,
$
(\mathbf{u} + \mathbf{v}) \cdot \mathbf{r}
=
\mathbf{v} \cdot \mathbf{s}_{\mathbf{r}}
$
for all $\mathbf{r} \in \operatorname{res}(\mathcal{C})$.
Thus, $(\mathbf{u} \, | \, \mathbf{v}) \in \Phi(\mathcal{C})^\perp$ if and only if
$\mathbf{v} \in \operatorname{tor}(\mathcal{C})^\perp$ and
$
(\mathbf{u} + \mathbf{v}) \cdot \mathbf{r}
=
\mathbf{v} \cdot \mathbf{s}_{\mathbf{r}}
$
for all  $\mathbf{r} \in \operatorname{res}(\mathcal{C})
$,
which yields the claimed description of $\Phi(\mathcal{C})^\perp$.

Finally, we check that this description is independent of the choice of
representatives $\mathbf{s}_{\mathbf{r}}$. Suppose we replace each
$\mathbf{s}_{\mathbf{r}}$ by
$
\mathbf{s}'_{\mathbf{r}} = \mathbf{s}_{\mathbf{r}} + \mathbf{t}_{\mathbf{r}}
$
where $\mathbf{t}_{\mathbf{r}} \in \operatorname{tor}(\mathcal{C})$.
For $\mathbf{v} \in \operatorname{tor}(\mathcal{C})^\perp$, we have
\begin{align*}
\mathbf{v} \cdot \mathbf{s}'_{\mathbf{r}}
=
\mathbf{v} \cdot \mathbf{s}_{\mathbf{r}}
+
\mathbf{v} \cdot \mathbf{t}_{\mathbf{r}}
=
\mathbf{v} \cdot \mathbf{s}_{\mathbf{r}},
\end{align*}
since $\mathbf{v} \cdot \mathbf{t}_{\mathbf{r}} = 0$ for all
$\mathbf{t}_{\mathbf{r}} \in \operatorname{tor}(\mathcal{C})$. Hence, the condition
$
(\mathbf{u} + \mathbf{v}) \cdot \mathbf{r}
=
\mathbf{v} \cdot \mathbf{s}_{\mathbf{r}}
$
for all $\mathbf{r} \in \operatorname{res}(\mathcal{C})$
is equivalent to
$
(\mathbf{u} + \mathbf{v}) \cdot \mathbf{r}
=
\mathbf{v} \cdot \mathbf{s}'_{\mathbf{r}}
$
for all $\mathbf{r} \in \operatorname{res}(\mathcal{C})
$,
so the resulting set $\Phi(\mathcal{C})^\perp$ does not depend on the chosen
system of representatives.
\end{proof}

\end{document}
